\chapter{Introduction}
\label{ch:intro}

\section{Motivation}

Large language models are expensive to run. Most of their energy goes into multiplying activations by weights
and into moving those weights from memory to the arithmetic units. Spiking neural networks (SNNs), inspired by
the brain, take a different approach: neurons exchange binary \emph{spikes}, so a synapse only does work when its
input neuron fires, and that work is an addition of the weight rather than a multiplication. If most neurons are
silent most of the time, a spiking network could need far less energy than a conventional one.

Transformers \citep{vaswani2017attention} dominate sequence modelling today, and several groups have built
spiking versions of them \citep{zhou2023spikformer, yao2023spikedriven, zhu2023spikegpt}. These models usually
report energy savings of an order of magnitude or more. Two questions motivated this thesis: how close can a
spiking transformer come to a conventional one on a real generation task, and how much energy does it actually
save once it is designed for hardware and compared fairly?

The task chosen is symbolic music generation. Folk tunes written in ABC notation are short, highly structured
texts: a tune has a meter, a mode, bars whose note durations must add up to the meter, and repeat signs that must
pair up. This makes it possible to judge a model not only by how well it predicts the next symbol but also by
whether the music it writes is well-formed, and the folk-rnn corpus \citep{sturm2016folkrnn} provides an
established LSTM baseline.

\section{Goals}

\begin{enumerate}
  \item Build a spiking transformer that generates folk tunes in ABC notation, and bring its quality as close as
        possible to a non-spiking transformer of the same size.
  \item Understand which design choices matter for spiking attention and training, including the ones that fail.
  \item Turn the spiking block into a datapath that hardware can implement with integer arithmetic only, and
        measure the accuracy cost.
  \item Estimate the energy per generated token and compare it fairly with conventional models, including
        8-bit quantized ones.
\end{enumerate}

\section{Contributions}

\begin{itemize}
  \item \textbf{\csq{} attention}, a spiking attention rule that weights each key by the square of its spike
        overlap with the query. It improves test bits/token by 0.10 over a learned count threshold, raises the
        share of correctly timed bars in generated tunes from 50\% to 58\% in short runs, and avoids an
        instability of integer thresholds that this thesis analyses.
  \item \textbf{A training recipe} for spiking transformers on a small dataset: distillation from a same-size
        ANN with a reverse attention loss, an attention-spread penalty, homeostasis, decaying schedules and spike
        dropout. The final model reaches 1.412 bits/token, ahead of a folk-rnn-style LSTM (1.591).
  \item \textbf{A multiplier-free integer datapath} for the spiking block (shift-coded residual inputs, int8
        weights, integer layer normalisation with a 16-entry reciprocal table, normalisation folded into 12-bit
        neurons), trained with quantization-aware training, with a bit-exact reference model and a hardware
        specification.
  \item \textbf{A fair energy comparison} against fp32 and int8 versions of the ANN and the LSTM, including memory
        traffic and a sensitivity analysis.
\end{itemize}

\section{Outline}

Chapter~\ref{ch:background} explains the necessary background: language modelling, transformers, LSTMs, spiking
neurons, distillation, quantization and ABC notation. Chapter~\ref{ch:related} reviews related work.
Chapter~\ref{ch:methods} describes the spiking model, \csq{} attention and the training recipe, and
Chapter~\ref{ch:hardware} the integer datapath and energy model. Chapter~\ref{ch:results} presents the
experiments: a preliminary phase on English text, the development of the folk-music model, the attention
variants, the step to integer arithmetic, an analysis of the final model, and the comparison with other models.
Chapter~\ref{ch:discussion} discusses what the results mean and their limitations, and
Chapter~\ref{ch:conclusion} concludes.
